\subsubsection*{Markman's candidate in dimension eight}

Above $n=3$ there is one concrete candidate of this kind. Markman proposes
it in \cite[Section 12]{Mar25b} and notes that it remains to be checked
against the weakened criterion. Its Chern character is a finite computation
in the ring $\QQ[\Theta]/(\Theta^{5})$. We carry the computation out: it
places the candidate at a rational point of the secant plane of
\Cref{def:secantplane} and reduces what has to be checked to a single
condition, which \Cref{thm:candidatefails} decides in the negative.

\begin{setupenv}[Markman's object]\label{setup:markman}
Let $(X,\Theta)$ be a principally polarised abelian fourfold, $d\ge3$ odd and
squarefree, $K=\QQ(\sqrt{-d})$ and $N=(d+9)/2$. Let $D_{i}$, $i\in\ZZ/N\ZZ$,
be $N$ translates of $\Theta$ in general position: any two, three or four of
them meet transversally, in a surface, a curve and $\Theta^{4}=24$ points,
and any five have empty intersection. Put $Z_{i}=D_{i}\cap D_{i+1}$,
$Z=\bigcup_{i}Z_{i}$, let $C_{i}=Z_{i}\cap Z_{i+1}=D_{i}\cap D_{i+1}\cap D_{i+2}$
be the curves along which consecutive components meet, and let
$\nu\colon\widetilde{Z}\to Z$ be the partial normalisation of $Z$ at
$m=(d+9)(2d-1)$ of its isolated points of self-intersection, the points of
$Z_{i}\cap Z_{j}$ for $i,j$ at cyclic distance at least three. Put
\[
  \cF=\bigl[\cO_{X}\longrightarrow\nu_{*}\cO_{\widetilde{Z}}\bigr]\otimes\cO_{X}(3\Theta),
\]
the complex placed in degrees $0$ and $1$.
\end{setupenv}

\begin{proposition}[Markman's candidate on the secant plane]
\label{prop:markman8}
In \Cref{setup:markman}, with $t=\Theta$ and $u_{t},v_{t}$ as in
\eqref{eq:uv} for $n=4$:
\begin{enumerate}[label=\textup{(\roman*)},nosep]
\item $\chi(\cO_{Z_{i}})=14$, $\chi(\cO_{C_{i}})=-36$, the isolated
  self-intersection points of $Z$ number $12N(N-5)=(d+9)(3d-3)$, and
  $\chi(\cO_{\widetilde Z})=110N-12N^{2}+2N(2d-1)=(d+9)(27-d)$;
\item $\ch(\nu_{*}\cO_{\widetilde Z})=N\Theta^{2}-2N\Theta^{3}
  +\chi(\cO_{\widetilde Z})\,[\mathrm{pt}]$;
\item $\ch(\cF)=u_{t}+3v_{t}$, that is
  \[
    \ch(\cF)=1+3\Theta-\tfrac{d}{2}\Theta^{2}-\tfrac{d}{2}\Theta^{3}
      +\tfrac{d^{2}}{24}\Theta^{4},
  \]
  the rational point $(a,b)=(1,3)$ of the secant plane $P_{\Theta}$; with
  the twist $\cO_{X}(\Theta)$ in place of $\cO_{X}(3\Theta)$ the degree two
  component would instead be $-\tfrac{d+8}{2}\Theta^{2}$, so the twist in
  \cite[Section 12]{Mar25b} must be read as $3\Theta$;
\item $\Nm_{K/\QQ}(3+\sqrt{-d})=d+9=2N$, with $N$ the number of divisors, and by
  \Cref{thm:secantinvariants} the Bogomolov discriminant of $\cF$ is
  $\Delta(\cF)=(d+9)\Theta^{2}$;
\item the objects $\Phi(\cF\boxtimes\cF)$ and $\Phi(\cF\boxtimes\cF^{\vee})$
  on $X\times\widehat X$ have ranks
  $\int_{X}\ch(\cF)^{2}=8d(d-9)$ and
  $\int_{X}\ch(\cF)\ch(\cF^{\vee})=8d(d+9)$, up to sign, both nonzero for
  squarefree $d$; the same formula gives $-8q$ for the pair of secant sheaves
  of \cite[Section 11]{Mar25b} on a threefold, which is his rank $8q$ up to
  the shift.
\end{enumerate}
Consequently $\cF\boxtimes\cF$ satisfies conditions (1), (2b) and (2c) of
Markman's strategy \cite[Section 4]{Mar25b}: the transform has nonzero rank,
its normalised Chern character is invariant under the special Mumford--Tate
group by \cite[Corollary 9.2]{Mar25b}, and its Weil component is nonzero by
\cite[Proposition 10.2]{Mar25b}, the relevant coefficient being
$\Nm(c)$ with $c=\tfrac12(1+3/\sqrt{-d})\ne0$. What remains is condition
(2a) in its weakened form: for $E=\Phi(\cF\boxtimes\cF)$, the map
$\operatorname{ev}_{E}$ must kill every class in
$H^{1}(T)\oplus H^{0}(\bigwedge^{2}T)$ that preserves $\kappa(E)$ to first
order.
\end{proposition}

\begin{proof}
On a principally polarised abelian fourfold $\chi(\cO_{X}(k\Theta))=k^{4}$
and $\Theta^{4}=24$. The Koszul complex of $r$ transversal translates of
$\Theta$ gives $\ch$ of the structure sheaf of their intersection as
$(1-e^{-\Theta})^{r}$, so $\ch(\cO_{Z_{i}})=\Theta^{2}-\Theta^{3}+\tfrac{7}{12}\Theta^{4}$,
$\ch(\cO_{C_{i}})=\Theta^{3}-\tfrac32\Theta^{4}$ and $\ch$ of four divisors
is $\Theta^{4}$. Integrating gives $\chi(\cO_{Z_{i}})=14$,
$\chi(\cO_{C_{i}})=-36$ and $24$ points; the same numbers are the Koszul sums
$0-2\cdot1+16$, $0-3+48-81$ and $0-4+96-324+256$ of the values
$\chi(\cO_{X}(-k\Theta))=k^{4}$.

(i), (ii). At every point $Z$ is locally a union of coordinate subspaces of
$\CC^{4}$: along $C_{i}$ two surfaces $\{x_{1}=x_{2}=0\}\cup\{x_{2}=x_{3}=0\}$,
at a point of $C_{i}\cap C_{i+1}$ three of them, at an isolated point
$\{x_{1}=x_{2}=0\}\cup\{x_{3}=x_{4}=0\}$. The ideals of such an arrangement
are monomial and distribute over sums, so the Mayer--Vietoris complex
\[
  0\to\cO_{Z}\to\bigoplus_{i}\cO_{Z_{i}}\to\bigoplus_{i<j}\cO_{Z_{i}\cap Z_{j}}
  \to\bigoplus_{i<j<k}\cO_{Z_{i}\cap Z_{j}\cap Z_{k}}\to\cdots
\]
with scheme-theoretic intersections is exact. Pairs at cyclic distance one
meet in $C_{i}$, $N$ of them; pairs at distance at least two meet in the $24$
points of four divisors, $N(N-3)/2$ of them; the only triples with nonempty
intersection are the consecutive ones $Z_{i}\cap Z_{i+1}\cap Z_{i+2}=C_{i}\cap C_{i+1}$,
$24$ points each, $N$ of them; and no four components meet, since that would
involve at least five divisors. Hence
\[
  \ch(\cO_{Z})=N\ch(\cO_{Z_{i}})-N\ch(\cO_{C_{i}})-\tfrac{N(N-3)}{2}\Theta^{4}+N\Theta^{4},
\]
whose degree two part is $N\Theta^{2}$, degree three part
$-N\Theta^{3}-N\Theta^{3}=-2N\Theta^{3}$, and degree four part
$\chi(\cO_{Z})[\mathrm{pt}]$ with $\chi(\cO_{Z})=14N+36N-12N(N-3)+24N=110N-12N^{2}$.
The isolated singular points are those of pairs at distance at least three,
of which there are $N(N-5)/2$, giving $12N(N-5)=3(d+9)(d-1)$ points; the
pairs at distance two lie on the double curves and are not isolated. At an
isolated point the two branches meet transversally, and separating them
gives $0\to\cO_{Z}\to\nu_{*}\cO_{\widetilde Z}\to k(p)\to0$, so each of the
$m$ normalisations raises $\chi$ by one and leaves the lower degrees of
$\ch$ unchanged. With $N=(d+9)/2$,
$\chi(\cO_{\widetilde Z})=110N-12N^{2}+2N(2d-1)=(d+9)(55-3(d+9)+2d-1)=(d+9)(27-d)$.

(iii) $\ch(\cF)=e^{3\Theta}\bigl(1-\ch(\nu_{*}\cO_{\widetilde Z})\bigr)$. The
degree two part is $(\tfrac92-N)\Theta^{2}=-\tfrac d2\Theta^{2}$, which is
where $N=(d+9)/2$ enters; with the twist $k\Theta$ it would be
$(\tfrac{k^{2}}{2}-N)\Theta^{2}$, equal to $-\tfrac d2\Theta^{2}$ exactly for
$k^{2}=9$, and $k=-3$ is excluded by the degree three part. The degree three
part is $(\tfrac{27}{6}-3N+2N)\Theta^{3}=(\tfrac92-N)\Theta^{3}=-\tfrac d2\Theta^{3}$,
and the degree four part is
$(\tfrac{81}{24}-\tfrac{9N}{2}+6N-\tfrac{\chi(\cO_{\widetilde Z})}{24})\Theta^{4}
=\tfrac{81+36N-\chi(\cO_{\widetilde Z})}{24}\Theta^{4}=\tfrac{d^{2}}{24}\Theta^{4}$
by (i). Comparing with \eqref{eq:uv} at $n=4$,
$u_{t}=1-\tfrac d2t^{2}+\tfrac{d^{2}}{24}t^{4}$ and $v_{t}=t-\tfrac d6t^{3}$,
gives $\ch(\cF)=u_{t}+3v_{t}$.

(iv) This follows from $\Nm(3+\sqrt{-d})=9+d$ and \eqref{eq:discnorm} with
$(a,b)=(1,3)$.

(v) Orlov's equivalence is $\Phi=(\id\times\Phi_{\cP})^{-1}\circ\mu^{*}$ with
$\mu(x,y)=(x+y,y)$, so the restriction of $\Phi(\cF_{1}\boxtimes\cF_{2})$ to
$\{x\}\times\widehat X$ is the Fourier--Mukai transform of
$\tau_{x}^{*}\cF_{1}\otimes\cF_{2}$, whose derived fibre at $L\in\widehat X$
is $R\Gamma(X,\tau_{x}^{*}\cF_{1}\otimes\cF_{2}\otimes L^{\pm1})$; the rank
of a perfect complex is the Euler characteristic of its derived fibre, and
$\chi(\tau_{x}^{*}\cF_{1}\otimes\cF_{2}\otimes L)=\int_{X}\ch(\cF_{1})\ch(\cF_{2})$
by Riemann--Roch, translations and $L\in\Pic^{0}$ acting trivially on
cohomology. With $\ch(\cF^{\vee})=u_{t}-3v_{t}$, the involution $\tau$
changing the sign of the odd degrees, $\int u_{t}^{2}=8d^{2}$,
$\int u_{t}v_{t}=0$ and $\int v_{t}^{2}=-8d$ give the two values. On a
threefold $\int(u_{t}+v_{t})^{2}=2\int u_{t}v_{t}=-8q$.

For the final assertion, condition (2b) is the
$\mathrm{Spin}(V_{\QQ})_{B}$-invariance of $\kappa(E)$, which
\cite[Corollary 9.2]{Mar25b} derives from $\ch(\cF)\in B$ and the
nonvanishing of the rank. Condition (2c) is
\cite[Proposition 10.2]{Mar25b}, which for $K$ imaginary quadratic reduces
to the nonvanishing of the component of $\ch(\cF)\otimes\ch(\cF)$ in
$\ell_{W}\otimes\ell_{\bbar W}$, and writing $\ch(\cF)=c\,\ell_{W}+\bbar c\,\ell_{\bbar W}$
with $\ell_{W}=e^{\sqrt{-d}\Theta}$ gives that coefficient as $c\bbar c$.
\end{proof}

\begin{remark}[The remaining condition]\label{rem:onecondition}
\Cref{prop:markman8} does not verify Markman's candidate. It verifies every
statement about the candidate that concerns classes, and it shows that the
one remaining statement concerns the object. For
$E=\Phi(\cF\boxtimes\cF)$ on the eightfold $X\times\widehat X$ of Weil type,
with $\kappa(E)=\ch(E)e^{-c_{1}(E)/\operatorname{rk}E}$, that statement reads:
every $x\in H^{1}(T)\oplus H^{0}(\bigwedge^{2}T)$ with
$x\lrcorner\,\kappa(E)=0$ satisfies $\operatorname{ev}_{E}(x)=0$ in
$\Ext^{2}(E,E)$. The $H^{2}(\cO)$ summand imposes nothing, because
$\operatorname{rk}E\ne0$ makes $z\cdot\kappa(E)\ne0$ for $z\ne0$. By
\Cref{thm:linebundle} the condition fails as soon as $E$ has a rank one
direct summand whose class is not a multiple of the polarisation, so the
first question about $E$ is whether it is indecomposable. Markman's object
at $n=3$ is reflexive of rank $8q$ away from a divisor, and he proves the
condition by descending to a quotient on which $\sigma_{E}$ is injective on
an invariant subspace \cite[Proposition 11.5]{Mar25b}. Whether the same can
be done for $\Phi(\cF\boxtimes\cF)$ is a question about the Ext groups of one
object on one eightfold, not about any Hodge class. \Cref{thm:factor} below
moves the question from the eightfold to the fourfold and from $E$ to $\cF$
itself: the condition is equivalent to the single identity
\eqref{eq:heisenberg} in $\Ext^{2}(\cF,\cF\otimes\Omega^{2}_{X})$, the other
half of it holding by construction. \Cref{thm:candidatefails} shows that the
identity fails, so the answer is negative.
\end{remark}

\subsubsection*{Reduction to the fourfold}

The condition of \Cref{rem:onecondition} concerns an object on an eightfold.
It is equivalent to a condition on the factor $\cF$ on the fourfold, and
there it splits into two conditions, one of which holds for Markman's object
because the object is constructed uniformly over the moduli of principally
polarised fourfolds. The reduction is not special to that object. It works
for every secant object on every abelian $n$-fold with $n\ge3$, and its
ingredients are three computations in the Hodge algebra of a single abelian
variety, which we carry out by hand. We first fix
conventions.

\begin{setupenv}[The weakened criterion for one object]\label{setup:criterion}
Let $Y$ be an abelian variety and $F$ a nonzero perfect complex on $Y$. Write
$HT^{2}(Y)=H^{0,2}(Y)\oplus H^{1}(T_{Y})\oplus H^{0}(\bigwedge^{2}T_{Y})$ as
in \eqref{eq:HT2}, with $x=(z,v,\pi)$ acting on $H^{*}(Y,\CC)$ by
$x\lrcorner$, and likewise $HT^{1}(Y)=H^{0,1}(Y)\oplus H^{0}(T_{Y})$, with
$(z,\theta)$ acting by $z\wedge$ and by the interior product $\theta\lrcorner$.
Normalise the Atiyah class $\mathrm{At}(F)\in\Ext^{1}(F,F\otimes\Omega^{1}_{Y})$
so that $\ch(F)=\operatorname{tr}\exp\mathrm{At}(F)$; for a line bundle $L$
it is then $c_{1}(L)$, as in \eqref{eq:obdef}.%
\footnote{This is the convention of \cite{Tod09} and of
\cite[footnote 14]{Mar25a}. Buchweitz and Flenner use the opposite sign,
which replaces $\mathrm{At}(F)$ by $-\mathrm{At}(F)$; this leaves
$\mathrm{At}(F)^{2}$ unchanged and changes only the sign of the middle term
of $\operatorname{ev}_{F}$, which does not affect the vanishing statements
below.}
The evaluation map of \Cref{thm:weakfails} is, under HKR,
\begin{equation}\label{eq:evgeneral}
  \operatorname{ev}_{F}(z,v,\pi)=z\cdot\id_{F}+v\lrcorner\,\mathrm{At}(F)
  +\tfrac12\,\pi\lrcorner\,\mathrm{At}(F)^{2}\;\in\;\Ext^{2}(F,F),
\end{equation}
the contraction of the polyvector with the exponential Atiyah class
\cite[Theorem A]{Hua21}, \cite{BF08}; its restriction to $H^{1}(T_{Y})$ is the
obstruction map, whose kernel is the space of first order deformations of $Y$
to which $F$ extends \cite{Tod09}, \cite[Section 8.3]{Mar25a}. Put
\[
  \Ann_{HT^{2}}(\ch F)=\bigl\{x\in HT^{2}(Y):x\lrcorner\,\ch(F)=0\bigr\},
\]
the Hochschild classes preserving $\ch(F)$ to first order. We say that $F$
\emph{satisfies the weakened criterion} if $\operatorname{ev}_{F}$ vanishes on
$\Ann_{HT^{2}}(\ch F)$. By the commutative square
$\sigma_{F}\circ\operatorname{ev}_{F}=(\,\cdot\,)\lrcorner\,\ch(F)$ of
\cite[Proposition 6.2.1 and Corollary 6.3.2]{BF08}, in the form quoted in
\cite[Lemma 11.3]{Mar25b},
\[
  \operatorname{ev}_{F}\bigl(\Ann_{HT^{2}}(\ch F)\bigr)
  =\ker\sigma_{F}\cap\image\operatorname{ev}_{F},
\]
so the weakened criterion is exactly the hypothesis of
\cite[Question 11.4]{Mar25b}, that the restriction of $\sigma_{F}$ to the
image of $\operatorname{ev}_{F}$ is injective. For a split object this is
\Cref{prop:evaluation}(ii).
\end{setupenv}

\begin{lemma}[The $B$-field transform]\label{lem:bfield}
Let $B\in H^{1,1}(Y)$. For $x=(z,v,\pi)\in HT^{2}(Y)$ put
\begin{equation}\label{eq:bfield}
  e^{B}x=\bigl(z+v\lrcorner\,B+\tfrac12\,\pi\lrcorner\,B^{2},\;
  v+\pi\lrcorner\,B,\;\pi\bigr),
\end{equation}
where $\pi\lrcorner\,B\in\Hom(H^{1,0},H^{0,1})=H^{1}(T_{Y})$ is the contraction
of one slot of $\pi$ with $B$: for $\pi=\partial_{a}\wedge\partial_{b}$ it sends
$p_{b}$ to $\partial_{a}\lrcorner\,B$, $p_{a}$ to $-\partial_{b}\lrcorner\,B$
and the other basis vectors to zero. Then $e^{B}$ is a linear automorphism of
$HT^{2}(Y)$ with inverse $e^{-B}$, and:
\begin{enumerate}[label=\textup{(\roman*)},nosep]
\item $x\lrcorner\,(w\,e^{B})=\bigl((e^{B}x)\lrcorner\,w\bigr)\,e^{B}$ for
  every even class $w\in H^{\mathrm{ev}}(Y,\CC)$;
\item for a line bundle $L$ with $c_{1}(L)=B$,
  $\operatorname{ev}_{F\otimes L}(x)=\operatorname{ev}_{F}(e^{B}x)$ and
  $\Ann_{HT^{2}}(\ch(F\otimes L))=e^{-B}\Ann_{HT^{2}}(\ch F)$; hence $F$
  satisfies the weakened criterion if and only if $F\otimes L$ does.
\end{enumerate}
The computation in (ii) is formal: it uses only that the twisted object has
Atiyah class $\mathrm{At}(F)+B\cdot\id_{F}$ and Chern character
$\ch(F)e^{B}$. It therefore applies verbatim to the twisted sheaves of
\cite[Section 2]{Mar25b}, with $B$ rational. In particular, for $F$ of
nonzero rank $r$, the criterion for the twisted object $F\otimes L^{-1}$ with class
$\kappa(F)=\ch(F)e^{-c_{1}(F)/r}$ is the same condition as the criterion for
$F$ with class $\ch(F)$.
\end{lemma}

\begin{proof}
The computation takes place in a graded algebra on which the interior
products $\iota_{a}=\partial_{a}\lrcorner$ act as odd derivations. It uses
only that $w$ is even; that $B$ is even and central with
$\iota_{a}\iota_{b}B=0$, which holds for $B$ of type $(1,1)$ and
$\partial_{a},\partial_{b}$ of type $(1,0)$; and that the odd scalar classes
$\iota_{a}B$ anticommute with every odd element. In $H^{*}(Y,\CC)$ these
properties are clear. They also hold in
$\bigoplus_{p,q}\Ext^{p}(F,F\otimes\Omega^{q}_{Y})=\Ext^{*}(F,F)\otimes
\bigwedge^{*}H^{1,0}$, on which the polyvectors act through the second factor
by odd derivations, because the classes coming from $H^{*}(Y,\CC)$ are graded
central in the Yoneda algebra. First, $\iota_{a}(B^{k})=k(\iota_{a}B)B^{k-1}$,
since $\iota_{a}B$ is odd and $B$ is even, so
$\iota_{a}e^{B}=(\iota_{a}B)e^{B}$. For $w$ even this gives
$v\lrcorner\,(we^{B})=(v\lrcorner\,w)e^{B}+w\sum_{a}v(p_{a})(\iota_{a}B)e^{B}
=\bigl(v\lrcorner\,w+(v\lrcorner\,B)w\bigr)e^{B}$, and
$z(we^{B})=(zw)e^{B}$. For the bivector, with $\pi\lrcorner=\iota_{a}\iota_{b}$,
\[
  \iota_{b}(we^{B})=(\iota_{b}w)e^{B}+w(\iota_{b}B)e^{B},
\]
and applying $\iota_{a}$, using that $\iota_{b}w$ and $\iota_{b}B$ are odd,
\[
  \iota_{a}\iota_{b}(we^{B})
  =\bigl[\iota_{a}\iota_{b}w
   +(\iota_{a}B)\iota_{b}w-(\iota_{b}B)\iota_{a}w
   +w\bigl((\iota_{a}\iota_{b}B)-(\iota_{b}B)(\iota_{a}B)\bigr)\bigr]e^{B}.
\]
The middle two terms are $(\pi\lrcorner\,B)\lrcorner\,w$ for the element
$\pi\lrcorner\,B$ of $\Hom(H^{1,0},H^{0,1})$ displayed in the statement, acting
as a derivation, and the last bracket is
$\tfrac12\iota_{a}\iota_{b}(B^{2})=\iota_{a}\bigl((\iota_{b}B)B\bigr)$.
Summing the three actions gives (i). Applying (i) with $B$ replaced by $-B$
and then by $B$, and using
$(\pi\lrcorner\,B)\lrcorner\,B=2(\iota_{a}B)(\iota_{b}B)=\pi\lrcorner\,B^{2}$
when $\iota_{a}\iota_{b}B=0$, shows $e^{-B}e^{B}=\id$.

(ii) The Atiyah class of $F\otimes L$ is $\mathrm{At}(F)+B\cdot\id_{F}$,
the Atiyah class being additive on tensor products \cite[Section 3]{BF03},
and the two summands commute, so
$\exp\mathrm{At}(F\otimes L)=\exp\mathrm{At}(F)\cdot e^{B}$. By (i) in the
algebra $\Ext^{*}(F,F)\otimes\bigwedge^{*}H^{1,0}$, with $w=\exp\mathrm{At}(F)$,
\[
  x\lrcorner\,\exp\mathrm{At}(F\otimes L)
  =\bigl((e^{B}x)\lrcorner\,\exp\mathrm{At}(F)\bigr)e^{B},
\]
and the right hand side has the same component in
$\Ext^{2}(F,F)\otimes\bigwedge^{0}$ as $(e^{B}x)\lrcorner\,\exp\mathrm{At}(F)$,
because $e^{B}-1=B+\tfrac12B^{2}+\cdots$ raises the exterior degree. This is
$\operatorname{ev}_{F\otimes L}(x)=\operatorname{ev}_{F}(e^{B}x)$. By (i) in
$H^{*}(Y,\CC)$ with $w=\ch(F)$, $x\lrcorner\,\ch(F\otimes L)=0$ if and only
if $(e^{B}x)\lrcorner\,\ch(F)=0$, which is the statement on the annihilators.
Hence $\operatorname{ev}_{F\otimes L}(\Ann_{HT^{2}}(\ch(F\otimes L)))
=\operatorname{ev}_{F}(\Ann_{HT^{2}}(\ch F))$.
\end{proof}

\begin{theorem}[Reduction to the factor]\label{thm:factor}
Let $(X,t)$ be a polarised abelian variety of dimension $n\ge3$, let $d\ge1$,
and let $F$ be a perfect complex on $X$ with
$\ch(F)=a\,u_{t}+b\,v_{t}$ for some $(a,b)\in\QQ^{2}\setminus\{0\}$, in the
notation of \eqref{eq:uv}. Then:
\begin{enumerate}[label=\textup{(\roman*)},nosep]
\item No nonzero class in $HT^{1}(X)$ preserves $\ch(F)$: the map
  $HT^{1}(X)\to H^{*}(X,\CC)$, $y\mapsto y\lrcorner\,\ch(F)$, is injective. Its
  two lowest components form the matrix
  $\bigl(\begin{smallmatrix}a&b\\ b&-ad\end{smallmatrix}\bigr)$ of
  determinant $-(a^{2}d+b^{2})=-\Nm_{K/\QQ}(b+a\sqrt{-d})\ne0$.
\item The classes in $HT^{2}(X)$ preserving $\ch(F)$ to first order are
  \begin{equation}\label{eq:secantkernel}
    \Ann_{HT^{2}}(\ch F)
    =\bigl\{(0,v,0):v\lrcorner\,t=0\bigr\}\;\oplus\;
     \bigl\{(\tfrac d2\,\pi\lrcorner\,t^{2},\,0,\,\pi):
       \pi\in H^{0}(\textstyle\bigwedge^{2}T_{X})\bigr\},
  \end{equation}
  the polarised first order deformations of $(X,t)$, of dimension $n(n+1)/2$,
  and the Poisson directions each compensated by a class in $H^{0,2}$, of
  dimension $n(n-1)/2$; the total is $n^{2}=\dim\cD_{n,n}$, and the space
  does not depend on $(a,b)$.
\item $F$ satisfies the weakened criterion if and only if
  \begin{enumerate}[label=\textup{(\alph*)},nosep]
  \item $v\lrcorner\,\mathrm{At}(F)=0$ in $\Ext^{2}(F,F)$ for every
    $v\in H^{1}(T_{X})$ with $v\lrcorner\,t=0$, that is, $F$ extends to first
    order along every polarised deformation of $(X,t)$; and
  \item $\mathrm{At}(F)^{2}=-d\,t^{2}\cdot\id_{F}$ in
    $\Ext^{2}(F,F\otimes\Omega^{2}_{X})$.
  \end{enumerate}
  The trace of (b) is automatic: $\operatorname{tr}\mathrm{At}(F)^{2}
  =2\ch_{2}(F)=-ad\,t^{2}=\operatorname{tr}(-d\,t^{2}\cdot\id_{F})$, so (b)
  is the vanishing of the trace free part of the square of the Atiyah class.
\item Let $F_{1},F_{2}$ be two such complexes, $G=F_{1}\boxtimes F_{2}$ on
  $X\times X$ and $E=\Phi(G)$ on $X\times\widehat X$, with $\Phi$ Orlov's
  equivalence as in \Cref{prop:markman8}(v). Then
  $\Ann_{HT^{2}}(\ch G)=\Ann_{HT^{2}}(\ch F_{1})\otimes1\oplus
  1\otimes\Ann_{HT^{2}}(\ch F_{2})$, of dimension $2n^{2}$, and
  \[
    E\ \text{satisfies the weakened criterion}
    \iff G\ \text{does}
    \iff F_{1}\ \text{and}\ F_{2}\ \text{do}.
  \]
\end{enumerate}
\end{theorem}

\begin{proof}
Choose a basis $p_{1},\dots,p_{n}$ of $H^{1,0}(X)$ and a basis
$q_{1},\dots,q_{n}$ of $H^{0,1}(X)$ with $t=\sum_{a}p_{a}\wedge q_{a}$, which
is possible because $t$ is a positive class of type $(1,1)$; let
$\partial_{a}$ be the dual basis of $T$, so that $\iota_{a}p_{b}=\delta_{ab}$
and $\iota_{a}q_{b}=0$. Write $\ch(F)=\sum_{k}c_{k}t^{k}$ with
$c_{k}\,k!=a(-d)^{k/2}$ for $k$ even and $c_{k}\,k!=b(-d)^{(k-1)/2}$ for $k$
odd, so that
\begin{equation}\label{eq:crecursion}
  c_{k+1}\,(k+1)!=-d\,c_{k-1}\,(k-1)!\qquad(k\ge1).
\end{equation}
Two facts about the Hodge algebra are used throughout: for $v\in H^{1}(T)$
acting as a derivation, $v\lrcorner\,t^{k}=k\,(v\lrcorner\,t)\,t^{k-1}$; and
for a bivector $\pi$ of type $(2,0)$ in the tangent directions,
$\pi\lrcorner\,t=0$ because $t$ has a single $(1,0)$ slot in each term, so
$\pi\lrcorner\,t^{k}=\binom{k}{2}(\pi\lrcorner\,t^{2})\,t^{k-2}$ by the footnote
to \eqref{eq:HT2}. Finally, multiplication by $t$ is injective on $H^{0,1}$
when $n\ge2$ and on $H^{0,2}$ when $n\ge3$: $q_{a}\wedge t$ contains the
monomial $p_{c}q_{c}q_{a}$ for any $c\ne a$, and $q_{a}q_{b}\wedge t$ contains
$p_{c}q_{c}q_{a}q_{b}$ for any $c\notin\{a,b\}$, and distinct monomials are
independent.

(i) For $y=(z,\theta)$ the class $y\lrcorner\,\ch(F)$ has components of type
$(k-1,k)$ equal to $c_{k-1}z\,t^{k-1}+c_{k}\,\theta\lrcorner\,t^{k}
=\bigl(c_{k-1}z+k\,c_{k}\,(\theta\lrcorner\,t)\bigr)t^{k-1}$. For $k=1$ this
is $az+b\,(\theta\lrcorner\,t)$ and for $k=2$ it is
$\bigl(bz-ad\,(\theta\lrcorner\,t)\bigr)t$, by \eqref{eq:crecursion}. Both
$z$ and $\theta\lrcorner\,t=\sum_{a}\theta^{a}q_{a}$ lie in $H^{0,1}$, on which
$t$ acts injectively, so $y\lrcorner\,\ch(F)=0$ forces
$(z,\theta\lrcorner\,t)$ into the kernel of the displayed matrix, whose
determinant is $-a^{2}d-b^{2}<0$. Hence $z=0$ and $\theta\lrcorner\,t=0$, and
the latter gives $\theta=0$ because $t$ is nondegenerate.

(ii) For $x=(z,v,\pi)$ the component of $x\lrcorner\,\ch(F)$ of type
$(k-1,k+1)$ is
\[
  c_{k-1}z\,t^{k-1}+c_{k}\,v\lrcorner\,t^{k}+c_{k+1}\,\pi\lrcorner\,t^{k+1}
  =\Bigl(c_{k-1}z+k\,c_{k}\,W+\tbinom{k+1}{2}c_{k+1}\,P\Bigr)t^{k-1},
\]
with $W=v\lrcorner\,t$ and $P=\pi\lrcorner\,t^{2}$, both in $H^{0,2}$. By
\eqref{eq:crecursion}, $\binom{k+1}{2}c_{k+1}=-\tfrac d2c_{k-1}$, so the bracket
is $c_{k-1}Z+k\,c_{k}W$ with $Z=z-\tfrac d2P$. For $k=1$ and $k=2$ this gives
$aZ+bW$ and $(bZ-adW)\,t$, and $t$ is injective on $H^{0,2}$ because $n\ge3$;
the same determinant forces $Z=W=0$. Conversely $Z=W=0$ kills every
component. So $\Ann_{HT^{2}}(\ch F)$ is the set of $(z,v,\pi)$ with
$v\lrcorner\,t=0$ and $z=\tfrac d2\,\pi\lrcorner\,t^{2}$, which is
\eqref{eq:secantkernel}; the two summands are independent because the first
has $\pi=0$ and the second is the graph of a linear map from the $\pi$
component. A $v\in\Hom(H^{1,0},H^{0,1})$ has $v\lrcorner\,t=\sum_{a}v(p_{a})\wedge q_{a}$,
which vanishes exactly when the matrix of $v$ is symmetric; this is the
tangent space of the polarised deformations, of dimension $n(n+1)/2$, as in
\Cref{prop:splitobstruction}(i). The bivectors contribute
$\dim\bigwedge^{2}T=n(n-1)/2$. Nothing in the description involves $(a,b)$.

(iii) By (ii) and \eqref{eq:evgeneral}, $\operatorname{ev}_{F}$ vanishes on
$\Ann_{HT^{2}}(\ch F)$ if and only if it vanishes on both summands. On the
first it is $v\mapsto v\lrcorner\,\mathrm{At}(F)$, which is (a). On the
second it is
\[
  \tfrac d2(\pi\lrcorner\,t^{2})\cdot\id_{F}+\tfrac12\,\pi\lrcorner\,\mathrm{At}(F)^{2}
  =\tfrac12\,\pi\lrcorner\bigl(\mathrm{At}(F)^{2}+d\,t^{2}\cdot\id_{F}\bigr),
\]
since $\pi\lrcorner\,(t^{2}\cdot\id_{F})=(\pi\lrcorner\,t^{2})\cdot\id_{F}$.
Now $\Omega^{2}_{X}=\cO_{X}\otimes\bigwedge^{2}H^{1,0}$, so
$\Ext^{2}(F,F\otimes\Omega^{2}_{X})=\Ext^{2}(F,F)\otimes\bigwedge^{2}H^{1,0}$,
and the contractions with all $\pi\in\bigwedge^{2}(H^{1,0})^{\vee}$ vanish on
an element of this space if and only if the element is zero, the pairing
being perfect. This is (b). For the trace, $\ch(F)=\operatorname{tr}\exp\mathrm{At}(F)$
gives $\ch_{2}(F)=\tfrac12\operatorname{tr}\mathrm{At}(F)^{2}$, and
$\ch_{2}(F)=c_{2}t^{2}=-\tfrac{ad}{2}t^{2}$, while
$\operatorname{tr}(\id_{F})=\rk F=c_{0}=a$.

(iv) The Hochschild cohomology of a product is the tensor product,
$HT^{2}(X\times X)=HT^{2}\otimes HT^{0}\oplus HT^{1}\otimes HT^{1}\oplus
HT^{0}\otimes HT^{2}$, the Atiyah class of $G$ is
$\mathrm{At}(F_{1})\boxtimes1+1\boxtimes\mathrm{At}(F_{2})$ \cite[Proposition 3.14]{BF03},
\cite[(8.4.1)]{Mar25a}, so $\exp\mathrm{At}(G)=\exp\mathrm{At}(F_{1})\boxtimes
\exp\mathrm{At}(F_{2})$, and $\Ext^{2}(G,G)$ is the corresponding K\"unneth
sum of $\Ext^{i}(F_{1},F_{1})\otimes\Ext^{j}(F_{2},F_{2})$. A polyvector on a
factor contracts only the forms of that factor, so for $x_{1}\in HT^{i}(X)$
and $x_{2}\in HT^{j}(X)$ with $i+j=2$,
\[
  (x_{1}\otimes x_{2})\lrcorner\,\ch(G)=\pm(x_{1}\lrcorner\,\ch F_{1})\otimes
  (x_{2}\lrcorner\,\ch F_{2}),\qquad
  \operatorname{ev}_{G}(x_{1}\otimes x_{2})=\pm\operatorname{ev}_{F_{1}}(x_{1})
  \otimes\operatorname{ev}_{F_{2}}(x_{2}),
\]
with $\operatorname{ev}(1)=\id$ on $HT^{0}$ and Koszul signs that play no
role; this is the decomposition \cite[(8.4.2), (8.4.3)]{Mar25a}. Let
$x=x_{20}+x_{11}+x_{02}$ be a class in $\Ann_{HT^{2}}(\ch G)$ according to the
three summands. Since $\ch(F_{i})$ is even, $x_{11}\lrcorner\,\ch(G)$ lies in
$H^{\mathrm{odd}}\otimes H^{\mathrm{odd}}$ while the other two terms lie in
$H^{\mathrm{ev}}\otimes H^{\mathrm{ev}}$, so $x_{11}\lrcorner\,\ch(G)=0$
separately; and $x_{11}\lrcorner\,\ch(G)$ is the image of $x_{11}$ under the
tensor product of the two maps of (i), which is injective, so $x_{11}=0$. For
the even part write $x_{20}=x_{1}\otimes1$ and $x_{02}=1\otimes x_{2}$; then
$(x_{1}\lrcorner\,\ch F_{1})\otimes\ch(F_{2})+\ch(F_{1})\otimes(x_{2}\lrcorner\,\ch F_{2})=0$.
Applying a linear form on the second factor taking the value $1$ on
$\ch(F_{2})$ shows $x_{1}\lrcorner\,\ch(F_{1})=\lambda\,\ch(F_{1})$ for a
scalar $\lambda$, and then $x_{2}\lrcorner\,\ch(F_{2})=-\lambda\,\ch(F_{2})$.
But every component of $x_{1}\lrcorner\,\ch(F_{1})$ has Hodge type $(p,p+2)$
while $\ch(F_{1})\ne0$ has type $(p,p)$; so $\lambda=0$ and
$x_{i}\in\Ann_{HT^{2}}(\ch F_{i})$. This proves the description of
$\Ann_{HT^{2}}(\ch G)$ and its dimension $2n^{2}$ by (ii). On it
$\operatorname{ev}_{G}(x_{1}\otimes1+1\otimes x_{2})=\pm\operatorname{ev}_{F_{1}}(x_{1})
\otimes\id_{F_{2}}\pm\id_{F_{1}}\otimes\operatorname{ev}_{F_{2}}(x_{2})$, whose
two terms lie in different K\"unneth summands and vanish exactly when
$\operatorname{ev}_{F_{i}}(x_{i})=0$, as $\id_{F_{i}}\ne0$. So $G$ satisfies
the criterion if and only if $F_{1}$ and $F_{2}$ do.

For the transport under $\Phi$ we use the commutative diagram of
\cite[Lemma 8.3.10]{Mar25a}: an equivalence $\Phi\colon D^{b}(A)\to D^{b}(B)$
of derived categories of abelian varieties induces isomorphisms
$\Phi^{HT}\colon HT^{2}(A)\to HT^{2}(B)$ and $\Phi^{H}\colon H^{*}(A,\CC)\to
H^{*}(B,\CC)$ with $\Phi^{H}(x\lrcorner\,\ch G)=\Phi^{HT}(x)\lrcorner\,\ch\Phi(G)$
\cite{CRV12}, and $\Phi\circ\operatorname{ev}_{G}=\operatorname{ev}_{\Phi(G)}\circ\Phi^{HT}$
on $\Hom(G,G[2])\to\Hom(\Phi G,\Phi G[2])$ \cite[Theorem A]{Hua21}; on abelian
varieties the Todd class is trivial, so there is no Duflo correction
\cite[Remark 8.3.11]{Mar25a}. The first square gives
$\Phi^{HT}\bigl(\Ann_{HT^{2}}(\ch G)\bigr)=\Ann_{HT^{2}}(\ch E)$ and the second
$\operatorname{ev}_{E}\bigl(\Ann_{HT^{2}}(\ch E)\bigr)
=\Phi\bigl(\operatorname{ev}_{G}(\Ann_{HT^{2}}(\ch G))\bigr)$, with $\Phi$ an
isomorphism on $\Ext^{2}$. So $E$ satisfies the criterion if and only if $G$
does.
\end{proof}

\Cref{fig:factor} draws the K\"unneth stack and the two conditions at
$n=4$.

\begin{figure}[tbp]
\centering
  \includegraphics[width=\textwidth]{fig_factor}
\caption{\Cref{thm:factor} for $n=4$. Top: the K\"unneth stack of
$HT^{2}(X\times X)$, whose middle summand $HT^{1}\otimes HT^{1}$ meets the
kernel in zero by (i), so that the classes preserving $\ch(F_{1}\boxtimes F_{2})$
are two copies of the kernel on one factor; across Orlov's isomorphism,
$HT^{2}$ of the eightfold with the $2n^{2}$ classes preserving $\ch(E)$.
Bottom: $HT^{2}(X)$ as three slabs with heights proportional to their
dimensions, and inside them the $n^{2}$ classes of \eqref{eq:secantkernel}:
the polarised deformations in front of $H^{1}(T_{X})$, and the whole of
$H^{0}(\bigwedge^{2}T_{X})$, each bivector tilted into $H^{0,2}$ by its
compensation $\tfrac d2\,\pi\lrcorner t^{2}$; nothing in $H^{0,2}$ alone
preserves $\ch(F)$. On the right, the two conditions that the evaluation map
imposes on the two parts, (a) closed by the construction of $\cF$ and (b)
the identity \eqref{eq:heisenberg}.}
\label{fig:factor}
\end{figure}

\begin{corollary}[The remaining condition on the fourfold]\label{cor:markmancondition}
In \Cref{setup:markman} let $E=\Phi(\cF\boxtimes\cF)$. Then $E$ satisfies the
weakened criterion of \cite[Question 11.4]{Mar25b} if and only if
\begin{equation}\label{eq:heisenberg}
  \mathrm{At}(\cF)^{2}=-d\,\Theta^{2}\cdot\id_{\cF}
  \qquad\text{in }\Ext^{2}(\cF,\cF\otimes\Omega^{2}_{X}).
\end{equation}
Equivalently, writing $\mathrm{At}(\cF)=\sum_{a}A_{a}\otimes p_{a}$ with
$A_{a}\in\Ext^{1}(\cF,\cF)$ in a basis $p_{a}$ of $H^{1,0}(X)$, and
$\Theta=\sum_{a}p_{a}\wedge q_{a}$ with $q_{a}\in H^{0,1}(X)$, all the
commutators $A_{a}A_{b}-A_{b}A_{a}\in\Ext^{2}(\cF,\cF)$ are the scalars
$\pm2d\,(q_{a}\wedge q_{b})\cdot\id_{\cF}$ prescribed by $\Theta^{2}$, the
sign being fixed by the Koszul rule.
\end{corollary}

\begin{proof}
By \Cref{prop:markman8}(iii), $\ch(\cF)=u_{\Theta}+3v_{\Theta}$, so
\Cref{thm:factor} applies with $n=4$, $(a,b)=(1,3)$ and $F_{1}=F_{2}=\cF$, and
$E$ satisfies the criterion if and only if $\cF$ satisfies (a) and (b). It
remains to show that (a) holds. The data of \Cref{setup:markman} are a
principally polarised fourfold $(X,\Theta)$ and $N$ points $x_{i}\in X$ with
$D_{i}=\tau_{x_{i}}^{*}\Theta$. The conditions imposed on the $D_{i}$, that
the intersections of two, three and four of them be smooth of the expected
codimension and that five of them be disjoint, are open in the parameter
space of such data over the local universal deformation $S$ of $(X,\Theta)$,
by the openness of smoothness for a proper flat family and the properness of
the projection. The isolated points of $Z$, being transversal intersections
of four divisors, form a finite \'etale cover of the open set where the
conditions hold, so the $m$ points to be normalised can be followed locally;
the partial normalisation separates two smooth branches at each of them and
commutes with base change. The construction therefore produces a perfect
complex $\cF_{S}$ on the universal family over a neighbourhood of the base
point in $S$, whose fibre at the base point is $\cF$ and whose base has
Kodaira--Spencer image the whole space of polarised deformations, of
dimension $10=n(n+1)/2$. The obstruction $v\lrcorner\,\mathrm{At}(\cF)$ to
extending $\cF$ along $v$ vanishes for every $v$ along which $\cF$ extends
\cite{Tod09}, and this is (a).

For the reformulation, $\Omega^{1}_{X}=\cO_{X}\otimes H^{1,0}$ gives
$\mathrm{At}(\cF)=\sum_{a}A_{a}\otimes p_{a}$, and
$\mathrm{At}(\cF)^{2}=-\sum_{a<b}(A_{a}A_{b}-A_{b}A_{a})\otimes p_{a}\wedge p_{b}$
by the Koszul rule, while $\Theta^{2}=-2\sum_{a<b}(p_{a}\wedge p_{b})(q_{a}\wedge q_{b})$
for $\Theta=\sum_{a}p_{a}\wedge q_{a}$. Comparing coefficients of
$p_{a}\wedge p_{b}$ gives the commutator relations.
\end{proof}

\begin{remark}[Scope of the reduction]\label{rem:factorclosed}
The reduction accomplishes three things. First, it removes the eightfold:
the condition of \Cref{rem:onecondition}, on an object of rank $8d(d-9)$ on
$X\times\widehat X$ of Weil type, is equivalent to one identity
\eqref{eq:heisenberg} for a rank one complex on the fourfold, in the Yoneda
algebra $\Ext^{*}(\cF,\cF)$. The dimension of the space of classes that has
to be killed drops from $2n^{2}$ to $n^{2}$ and, by (a), to $n(n-1)/2=6$. At
$n=3$ the count $n^{2}=9$ is the nine-dimensional kernel of
\cite[Corollaries 8.3.12 and 8.5.2]{Mar25a}, which \eqref{eq:secantkernel}
writes out for every secant object in every dimension. Second, the half of
that space made of geometric deformations imposes nothing, because Markman's
construction is uniform over the moduli of principally polarised fourfolds.
Third, by \Cref{lem:bfield} the two forms of the criterion, for $E$ with
$\ch(E)$ and for the twisted object with $\kappa(E)$, coincide, so nothing
is lost by working with the untwisted object.

What remains is the identity itself, which is a Heisenberg relation: the
components of the Atiyah class in the $n$ tangent directions must commute up
to the scalars prescribed by $\Theta^{2}$, and the Chern character forces
only the trace of the relation. The identity says exactly that $\cF$
deforms, as an object of the deformed category, along the Poisson directions
of $X$ compensated by the gerbe classes $\tfrac d2\,\pi\lrcorner\,\Theta^{2}$;
these are the noncommutative first order deformations of $X$ with a gerby
correction, in the language of \cite[Section 8.5]{Mar25a}. It holds in the
one case where the method is known to work. For Markman's sheaves
$\cF_{1}=I_{\bigcup C_{i}}(\Theta)$ on the Jacobian threefold of
\cite[Section 11.1]{Mar25b}, with $\ch=u_{\Theta}+v_{\Theta}$ and $d=q$, the
kernel of the obstruction map equals the annihilator of the Chern character
\cite[Corollary 8.3.12]{Mar25a}, which is stronger than the weakened
criterion; so by \Cref{thm:factor}(iii) these sheaves satisfy
$\mathrm{At}(\cF_{1})^{2}=-q\,\Theta^{2}\cdot\id$ on the threefold. For the
rank one complex of \Cref{setup:markman} on the fourfold, the identity is a
computation of Yoneda products in $\Ext^{*}(\cF,\cF)$, and the computation
can be made local. Below we carry it out at one point of the support and
find that the identity fails.
\end{remark}

\subsubsection*{Failure at an isolated double point}

The identity \eqref{eq:heisenberg} can be tested locally. The evaluation map
is an algebra homomorphism \cite[(8.3.3)]{Mar25a}, so on a bivector
$\theta\wedge\theta'$ it is the Yoneda product of the classes
\[
  A_{\theta}=\operatorname{ev}_{\cF}(\theta)=\theta\lrcorner\,\mathrm{At}(\cF)
  \in\Ext^{1}(\cF,\cF),\qquad\theta\in H^{0}(T_{X}),
\]
which are the first order deformations of $\cF$ along the translations of
$X$. The product of two global classes restricts to the product of their
germs. At a point where the support of $\cF$ has two smooth branches of
codimension two meeting transversally, the germs multiply nontrivially;
since the compensating classes $z\cdot\id_{\cF}$ have zero germs, the
identity fails. Markman's construction has such points and cannot avoid
them.

\begin{lemma}[Localisation]\label{lem:edge}
Let $F$ be a perfect complex on a smooth projective variety $Y$, let
$\cE xt^{m}(F,F)$ be the sheaf associated with $U\mapsto\Hom_{D(U)}(F|_{U},F|_{U}[m])$,
and let $e\colon\Ext^{m}(F,F)\to H^{0}(Y,\cE xt^{m}(F,F))$ send a morphism
$F\to F[m]$ to its germs. Then $e$ is compatible with Yoneda products;
$e(z\cdot\id_{F})=0$ for $z\in H^{m}(\cO_{Y})$, $m\ge1$; and for
$\theta\in H^{0}(T_{Y})$ the germ of $A_{\theta}$ at $p$ is the class of the
jet extension $0\to F_{p}\to J_{\theta}(F_{p})\to F_{p}\to0$ of the germ
$F_{p}$ along the derivation $\theta$ of $\cO_{Y,p}$, which is
$\cO_{Y,p}$-linear in $\theta$.
\end{lemma}

\begin{proof}
Composition of morphisms commutes with restriction to open sets; this is the
first assertion. For the second, $z\cdot\id_{F}$ restricted to an affine open
$U$ is $z|_{U}\cdot\id_{F|_{U}}$, and $z|_{U}\in H^{m}(U,\cO_{U})=0$. For the
third, the Atiyah class of $F$ is the class of the first jet sequence
$0\to F\otimes\Omega^{1}_{Y}\to J^{1}(F)\to F\to0$ \cite[Section 3]{BF03},
and its restriction to $U$ is the jet sequence of $F|_{U}$. Contracting with
$\theta$ is pushing out along $\id\otimes\theta\colon F\otimes\Omega^{1}\to F$,
which produces the extension $J_{\theta}(F)$ whose underlying sheaf is
$F\oplus F$ with $f\cdot(s,s')=(fs+\theta(f)s',fs')$. Linearity in $\theta$
is linearity of the contraction.
\end{proof}

\begin{lemma}[Two branches through a point]\label{lem:localproducts}
Let $\cO$ be the local ring of $\CC^{4}$ at the origin or its completion
$\CC[[x_{1},\dots,x_{4}]]$; let $P_{1}=\{x_{1}=x_{2}=0\}$ and
$P_{2}=\{x_{3}=x_{4}=0\}$, two smooth surfaces meeting transversally at the
origin; put $N_{1}=\langle\partial_{1},\partial_{2}\rangle$ and
$N_{2}=\langle\partial_{3},\partial_{4}\rangle$, the normal spaces of $P_{1}$
and of $P_{2}$ at the origin, so that
$\bigwedge^{2}T_{0}=\bigwedge^{2}N_{1}\oplus N_{1}\otimes N_{2}\oplus\bigwedge^{2}N_{2}$.
For a germ of vector field $\theta$ and a perfect complex $F$ of
$\cO$-modules write $a_{\theta}\in\Ext^{1}_{\cO}(F,F)$ for the class of the
jet extension along $\theta$.
\begin{enumerate}[label=\textup{(\roman*)},nosep]
\item For $F=I=I_{P_{1}}\cap I_{P_{2}}$, the ideal of the union,
  $\Ext^{2}_{\cO}(I,I)\cong N_{1}\otimes N_{2}$, a vector space of dimension
  four on which the maximal ideal acts trivially, and
  $a_{\theta}a_{\theta'}$ is, up to a fixed nonzero scalar, the component of
  $\theta(0)\wedge\theta'(0)$ in $N_{1}\otimes N_{2}$. In particular
  $a_{\theta}a_{\theta'}\ne0$ unless $\theta(0)\wedge\theta'(0)\in
  \bigwedge^{2}N_{1}\oplus\bigwedge^{2}N_{2}$.
\item For $F=K=[\cO\to\cO_{P_{1}}\oplus\cO_{P_{2}}]$, the complex in degrees
  $0$ and $1$ whose cohomology is $I$ in degree $0$ and the residue field $k$
  in degree $1$, $a_{\theta}a_{\theta'}\ne0$ unless
  $\theta(0)\wedge\theta'(0)\in N_{1}\otimes N_{2}$.
\end{enumerate}
In both cases the bilinear map $(\theta,\theta')\mapsto a_{\theta}a_{\theta'}$
is nonzero, and it depends only on $\theta(0)\wedge\theta'(0)$.
\end{lemma}

\begin{proof}
Write $A=\CC[x_{1},x_{2}]$, $B=\CC[x_{3},x_{4}]$, $R=A\otimes_{\CC}B$, with
maximal ideals $\mathfrak m_{A}$, $\mathfrak m_{B}$ at the origin, and
$k=\CC$. The ideal of the union is $I=\mathfrak m_{A}R\cap\mathfrak m_{B}R
=\mathfrak m_{A}\otimes_{\CC}\mathfrak m_{B}$, because the two ideals are
generated by disjoint sets of variables. Everything below is homogeneous, the
$\Ext$ groups in degree two are of finite length, and finite length modules
and jet classes are unchanged by the flat base change from $R$ to $\cO$; so we
compute over $R$. For finitely generated modules $M,M'$ over $A$ and $N,N'$
over $B$, the tensor product of free resolutions resolves $M\otimes_{\CC}N$
and the K\"unneth formula gives
\begin{equation}\label{eq:kunnethext}
  \Ext^{n}_{R}(M\otimes N,M'\otimes N')
  =\bigoplus_{i+j=n}\Ext^{i}_{A}(M,M')\otimes_{\CC}\Ext^{j}_{B}(N,N'),
\end{equation}
compatibly with Yoneda products up to the Koszul sign, since a product of
chain maps on the tensor product of resolutions is the tensor product of the
products.

The basic case is a single plane. The Koszul complex
$0\to A\xrightarrow{(-x_{2},x_{1})^{t}}A^{2}\xrightarrow{(x_{1},x_{2})}\mathfrak m_{A}\to0$
is a free resolution, so $\mathfrak m_{A}$ has projective dimension one and
$\Ext^{i}_{A}(\mathfrak m_{A},-)=0$ for $i\ge2$; and
$\Ext^{1}_{A}(\mathfrak m_{A},\mathfrak m_{A})=\coker\bigl(\Hom(A^{2},\mathfrak m_{A})
\to\Hom(A,\mathfrak m_{A})\bigr)=\mathfrak m_{A}/\mathfrak m_{A}^{2}$, a
two-dimensional space killed by $\mathfrak m_{A}$. On a complex of free
modules with differential $d$ the jet extension along a derivation $D$ is the
complex with terms $P_{i}\oplus P_{i}$ and differential
$\bigl(\begin{smallmatrix}d&-D(d)\\0&d\end{smallmatrix}\bigr)$, $D(d)$ the
entrywise derivative, because the twisted module $J_{D}(P_{i})$ is free on the
elements $(e,0)$ and $(0,e)$ and $d(0,e)=\sum_{l}d_{le}(0,e_{l})-\sum_{l}D(d_{le})(e_{l},0)$;
so the jet class is represented by the degree $-1$ chain map $-D(d)$. For
$D=\partial_{i}$ on the Koszul complex this is $e\mapsto\delta_{i2}f_{1}-\delta_{i1}f_{2}$
on the generator $e$ of the last term, whose image in
$\mathfrak m_{A}/\mathfrak m_{A}^{2}$ is $\delta_{i2}x_{1}-\delta_{i1}x_{2}$:
\[
  \alpha_{1}:=a_{\partial_{1}}=-[x_{2}],\qquad
  \alpha_{2}:=a_{\partial_{2}}=[x_{1}],
\]
a basis of $\Ext^{1}_{A}(\mathfrak m_{A},\mathfrak m_{A})$, which we identify
with $N_{1}$ by $\partial_{i}\mapsto\alpha_{i}$. Similarly $\beta_{3},\beta_{4}$
form a basis of $\Ext^{1}_{B}(\mathfrak m_{B},\mathfrak m_{B})$, identified
with $N_{2}$.

(i) By \eqref{eq:kunnethext} with $M=M'=\mathfrak m_{A}$, $N=N'=\mathfrak m_{B}$,
\[
  \Ext^{2}_{R}(I,I)=\Ext^{1}_{A}\otimes\Ext^{1}_{B}=N_{1}\otimes N_{2},
\]
the other two summands vanishing by the projective dimension, and it is killed
by $\mathfrak m_{A}$ and $\mathfrak m_{B}$. The derivation $\partial_{1}$ acts on
the factor $A$, so the jet extension of $I=\mathfrak m_{A}\otimes\mathfrak m_{B}$
along it is $J_{\partial_{1}}(\mathfrak m_{A})\otimes\mathfrak m_{B}$ and
$a_{\partial_{1}}=\alpha_{1}\otimes1$; likewise $a_{\partial_{2}}=\alpha_{2}\otimes1$,
$a_{\partial_{3}}=1\otimes\beta_{3}$, $a_{\partial_{4}}=1\otimes\beta_{4}$.
Hence $a_{\partial_{i}}a_{\partial_{j}}=\pm\alpha_{i}\otimes\beta_{j}$ for
$i\le2<j$, the four of them a basis of $N_{1}\otimes N_{2}$, while
$a_{\partial_{1}}a_{\partial_{2}}=\alpha_{1}\alpha_{2}\otimes1=0$ and
$a_{\partial_{3}}a_{\partial_{4}}=0$ because $\Ext^{2}_{A}(\mathfrak m_{A},\mathfrak m_{A})=0$.
For germs $\theta=\sum_{i}c^{i}\partial_{i}$ and $\theta'=\sum_{j}c'^{j}\partial_{j}$
the linearity of \Cref{lem:edge} and the triviality of the module structure
give $a_{\theta}a_{\theta'}=\sum_{i,j}c^{i}(0)c'^{j}(0)\,a_{\partial_{i}}a_{\partial_{j}}$,
the image of $\theta(0)\wedge\theta'(0)$ under the projection of
$\bigwedge^{2}T_{0}$ onto $N_{1}\otimes N_{2}$ followed by the isomorphism
$\partial_{i}\otimes\partial_{j}\mapsto\pm\alpha_{i}\otimes\beta_{j}$.

(ii) Let $Q=\cO_{P_{1}}\oplus\cO_{P_{2}}$. The map $\cO\to Q$ has kernel $I$
and cokernel $\cO/(I_{P_{1}}+I_{P_{2}})=k$, so $K$ sits in an exact triangle
$I\xrightarrow{\iota}K\xrightarrow{\pi}k[-1]\xrightarrow{\varepsilon}I[1]$.
The jet extension along $\theta$ is an exact functor, natural in its
argument, so the jet classes commute with $\iota$ and $\pi$; and for the
residue field $a^{k}_{\theta}$ depends only on $\theta(0)$, because
$\Ext^{1}(k,k)$ is killed by $\mathfrak m$ and the class is linear in
$\theta$. Therefore
\[
  \pi[2]\circ a^{K}_{\theta}a^{K}_{\theta'}
  =a^{k}_{\theta}a^{k}_{\theta'}\circ\pi
  =\pi^{*}\bigl(a^{k}_{\theta(0)}a^{k}_{\theta'(0)}\bigr),
  \qquad\pi^{*}\colon\Ext^{2}(k,k)\to\Ext^{1}(K,k),
\]
and it suffices to show that the right hand side is nonzero. The long exact
sequence of the triangle for $\Hom(-,k)$ reads
\[
  \Hom(I,k)\xrightarrow{\;\varepsilon^{*}\;}\Ext^{2}(k,k)\xrightarrow{\;\pi^{*}\;}
  \Ext^{1}(K,k),
\]
so $\ker\pi^{*}=\image\varepsilon^{*}$. Now $k=k_{A}\otimes k_{B}$ and, by
\eqref{eq:kunnethext}, $\Ext^{*}(k,k)$ is the tensor product of
$\Ext^{*}_{A}(k,k)$ and $\Ext^{*}_{B}(k,k)$. On the Koszul resolution
$A\to A^{2}\to A$ of $k_{A}$ the jet chain maps of $\partial_{1},\partial_{2}$
are $-D(d)$ as above, and their composite $A\to A$ is the scalar $-1$, whose
class in $\Ext^{2}_{A}(k,k)=\Hom(A,k)=k$ is nonzero; the degree one classes
$\gamma_{1},\gamma_{2}$ of $\partial_{1},\partial_{2}$ are the basis
$-e_{1}^{*},-e_{2}^{*}$ of $\Ext^{1}_{A}(k,k)=\Hom(A^{2},k)$, and their
product spans $\Ext^{2}_{A}(k,k)$. With the analogous classes
$\gamma_{3},\gamma_{4}$ over $B$, the products $\gamma_{1}\gamma_{2}\otimes1$,
$\gamma_{i}\otimes\gamma_{j}$ ($i\le2<j$) and $1\otimes\gamma_{3}\gamma_{4}$
form a basis of $\Ext^{2}(k,k)=k\oplus(k^{2}\otimes k^{2})\oplus k$, so
$\partial_{i}\wedge\partial_{j}\mapsto a^{k}_{\partial_{i}}a^{k}_{\partial_{j}}$
identifies $\bigwedge^{2}(N_{1}\oplus N_{2})$ with $\Ext^{2}(k,k)$, and
$a^{k}_{\theta(0)}a^{k}_{\theta'(0)}$ is the image of $\theta(0)\wedge\theta'(0)$. The
class $\varepsilon$ is the K\"unneth product $\varepsilon_{A}\otimes\varepsilon_{B}$
of the classes of $0\to\mathfrak m_{A}\to A\to k\to0$ and its $B$-analogue,
because $K$ is the two-step truncation of the tensor product of the complexes
$[A\to k_{A}]$ and $[B\to k_{B}]$ and $\varepsilon$ is the inclusion of the
omitted term $k_{A}\otimes k_{B}$. Hence $\varepsilon^{*}$ maps
$\Hom(I,k)=\Hom(\mathfrak m_{A},k)\otimes\Hom(\mathfrak m_{B},k)$ into
$\Ext^{1}_{A}(k,k)\otimes\Ext^{1}_{B}(k,k)=N_{1}\otimes N_{2}$, and onto it,
since $\Hom(\mathfrak m_{A},k)\to\Ext^{1}_{A}(k,k)$ is an isomorphism: it is
the connecting homomorphism of the sequence, injective because
$\Hom(A,k)\to\Hom(\mathfrak m_{A},k)$ is zero and surjective because
$\Ext^{1}_{A}(A,k)=0$. So $\ker\pi^{*}=N_{1}\otimes N_{2}$
and $\pi^{*}(\theta(0)\wedge\theta'(0))\ne0$ exactly when
$\theta(0)\wedge\theta'(0)\notin N_{1}\otimes N_{2}$.
\end{proof}

\begin{theorem}[Markman's candidate does not satisfy the weakened criterion]
\label{thm:candidatefails}
In \Cref{setup:markman}, for every odd squarefree $d\ge3$, every choice of the
translates $D_{i}$ in general position and every choice of the $m$ points to
be normalised:
\begin{enumerate}[label=\textup{(\roman*)},nosep]
\item the identity \eqref{eq:heisenberg} fails: there is a constant bivector
  $\pi\in H^{0}(\bigwedge^{2}T_{X})$ with
  $\tfrac12\pi\lrcorner\bigl(\mathrm{At}(\cF)^{2}+d\,\Theta^{2}\cdot\id_{\cF}\bigr)\ne0$
  in $\Ext^{2}(\cF,\cF)$;
\item $\cF$ does not satisfy the weakened criterion: the class
  $x_{\cF}=(\tfrac d2\,\pi\lrcorner\,\Theta^{2},0,\pi)$ has
  $x_{\cF}\lrcorner\,\ch(\cF)=0$ and $\operatorname{ev}_{\cF}(x_{\cF})\ne0$;
\item for every perfect complex $F_{2}$ on $X$ with Chern character on the
  secant plane, in particular for $F_{2}=\cF$ and $F_{2}=\cF^{\vee}$, the
  object $E=\Phi(\cF\boxtimes F_{2})$ on the eightfold $X\times\widehat X$ of
  Weil type does not satisfy the weakened criterion of
  \cite[Question 11.4]{Mar25b}: the class $x=\Phi^{HT}(x_{\cF}\otimes1)$
  satisfies $x\lrcorner\,\ch(E)=0$ and $\operatorname{ev}_{E}(x)\ne0$, so
  $\sigma_{E}$ is not injective on the image of $\operatorname{ev}_{E}$.
\end{enumerate}
\end{theorem}

\begin{proof}
(i) The surfaces $Z_{i}$ and $Z_{j}$ at cyclic distance at least three meet
in the $24$ points of $D_{i}\cap D_{i+1}\cap D_{j}\cap D_{j+1}$, and there are
$12N(N-5)=(d+9)(3d-3)$ such points, of which $m=(d+9)(2d-1)$ are normalised;
the remaining $(d+9)(d-2)\ge12$ are not. Let $p$ be one of the latter. The
four divisors through $p$ meet transversally, so their local equations
$x_{1},\dots,x_{4}$ form a regular system of parameters of $\cO_{X,p}$ and
identify its completion with $\CC[[x_{1},\dots,x_{4}]]$, with
$D_{i}=\{x_{1}=0\}$, $D_{i+1}=\{x_{2}=0\}$, $D_{j}=\{x_{3}=0\}$ and
$D_{j+1}=\{x_{4}=0\}$; hence $Z_{i}=P_{1}$ and $Z_{j}=P_{2}$ in the
notation of \Cref{lem:localproducts}. The completion is flat, so a class
whose completed germ is nonzero is nonzero. No other $D_{k}$ passes through
$p$, since five divisors have empty intersection, so no other component of
$Z$ does. Near $p$ the map $\nu$ is an isomorphism,
$\nu_{*}\cO_{\widetilde Z}=\cO_{Z}$, and $\cF$ is the ideal sheaf
$I_{Z}=I_{P_{1}}\cap I_{P_{2}}$ twisted by the locally trivial $\cO_{X}(3\Theta)$:
the germ $\cF_{p}$ is the module $I$ of \Cref{lem:localproducts}(i).

Let $\theta,\theta'\in H^{0}(T_{X})$ be constant vector fields whose values
at $p$ are $\partial_{1}$ and $\partial_{3}$; they exist because the constant
fields span every tangent space. Put $\pi=\theta\wedge\theta'$. Since
$\operatorname{ev}_{\cF}$ is an algebra homomorphism,
$\operatorname{ev}_{\cF}(\pi)=\tfrac12\pi\lrcorner\,\mathrm{At}(\cF)^{2}
=A_{\theta}A_{\theta'}$. By \Cref{lem:edge} its germ at $p$ is
$a_{\theta}a_{\theta'}$, which is nonzero by \Cref{lem:localproducts}(i),
because $\partial_{1}\wedge\partial_{3}$ has nonzero component in
$N_{1}\otimes N_{2}$. The germ of
$\tfrac d2(\pi\lrcorner\,\Theta^{2})\cdot\id_{\cF}$ at $p$ is zero by
\Cref{lem:edge}. So the germ at $p$ of
$\tfrac12\pi\lrcorner(\mathrm{At}(\cF)^{2}+d\Theta^{2}\cdot\id_{\cF})$ is
$a_{\theta}a_{\theta'}\ne0$, and the class itself is nonzero.

(ii) By \Cref{thm:factor}(ii) the class $x_{\cF}$ preserves $\ch(\cF)$, and by
\eqref{eq:evgeneral} $\operatorname{ev}_{\cF}(x_{\cF})$ is the class of (i).

(iii) Both $\cF$ and $\cF^{\vee}$ are admissible as $F_{2}$, since
$\ch(\cF^{\vee})=u_{\Theta}-3v_{\Theta}$ lies on the secant plane. By
\Cref{thm:factor}(iv), $E$ satisfies the criterion only if both factors do,
and $\cF$ does not. The class $x$ of the statement is the transport of
$x_{\cF}\otimes1\in\Ann_{HT^{2}}(\ch(\cF\boxtimes F_{2}))$, on which
$\operatorname{ev}$ is $\pm\operatorname{ev}_{\cF}(x_{\cF})\otimes\id_{F_{2}}\ne0$;
since $\Phi$ is an isomorphism on $\Ext^{2}$, $\operatorname{ev}_{E}(x)\ne0$.
The last assertion follows from the commutative square of
\Cref{setup:criterion}: $\sigma_{E}(\operatorname{ev}_{E}(x))=x\lrcorner\,\ch(E)=0$
with $\operatorname{ev}_{E}(x)\ne0$.
\end{proof}

\Cref{fig:doublepoint} draws the local computation at the double point.

\begin{remark}[Normalised points]\label{rem:normalisedpoints}
At a normalised point the germ of $\cF$ is the complex $K$ of
\Cref{lem:localproducts}(ii), and the same argument, with
$\theta(p)=\partial_{1}$ and $\theta'(p)=\partial_{2}$, produces a nonzero
germ. The failure therefore does not depend on which points are separated.
Separating all of them would not help even if the construction allowed it:
normalising every isolated point changes $\chi(\cO_{\widetilde Z})$ to
$110N-12N^{2}+12N(N-5)=50N$, while the secant plane requires
$(d+9)(27-d)=72N-4N^{2}$, and the two agree only for $N=11/2$. The two local
computations are complementary: at a point of type (i) the obstruction lies
in $N_{1}\otimes N_{2}$, at a point of type (ii) in
$\bigwedge^{2}N_{1}\oplus\bigwedge^{2}N_{2}$, and together these exhaust
$\bigwedge^{2}T_{p}$.
\end{remark}

\begin{figure}[tbp]
\centering
  \includegraphics[width=\textwidth]{fig_doublepoint}
\caption{\Cref{lem:localproducts} and \Cref{thm:candidatefails}. Left: two
smooth surfaces of a fourfold meeting transversally at $p$, with the normal
spaces $N_{1}$ and $N_{2}$, each tangent to the other branch. Centre: the
bivectors at $p$, split into $\bigwedge^{2}N_{1}$, $N_{1}\otimes N_{2}$ and
$\bigwedge^{2}N_{2}$, with heights proportional to the dimensions. Right: the
local $\Ext^{2}$ of the germ of $\cF$ at $p$ in the two cases, and the part
of the bivectors that the products of the translation classes carry onto it:
the mixed part when the point is not normalised, the two same-plane parts
when it is. The compensating scalar has no germ, so in either case the
identity \eqref{eq:heisenberg} fails at $p$.}
\label{fig:doublepoint}
\end{figure}

\begin{corollary}[A local obstruction in every dimension]\label{cor:localobstruction}
Let $(X,t)$ be a polarised abelian variety of dimension $n\ge4$ and $F$ a
perfect complex on $X$ with $\ch(F)=au_{t}+bv_{t}$, $(a,b)\ne(0,0)$. Suppose
that at some point $p$ the germ $F_{p}$ is isomorphic, in the derived
category of $\cO_{X,p}$-modules, to the ideal of the union of two smooth germs
of codimension two meeting transversally along a germ of codimension four,
or, when $n=4$, to the complex $[\cO\to\cO_{P_{1}}\oplus\cO_{P_{2}}]$ of two
such germs. Then $F$ does not satisfy the weakened criterion, and neither does
$\Phi(F\boxtimes F_{2})$ on $X\times\widehat X$ for any secant $F_{2}$. In
particular, in every dimension $n\ge4$, the union $Z$ of the intersections
$D_{i}\cap D_{i+1}$ of consecutive translates of a polarisation in general
position contains, for every pair $(i,j)$ of nonconsecutive indices, a
subvariety $Z_{i}\cap Z_{j}$ of codimension four along which $Z$ has two
transversal branches; any secant object of the form $[\cO_{X}\to\nu_{*}\cO_{\widetilde Z}]\otimes L$
in which $\nu$ is an isomorphism at a general point of one such
$Z_{i}\cap Z_{j}$ fails the criterion.
\end{corollary}

\begin{proof}
In formal coordinates $x_{1},\dots,x_{n}$ at $p$ with the two branches
$\{x_{1}=x_{2}=0\}$ and $\{x_{3}=x_{4}=0\}$, which exist because the four
local equations of the branches form part of a regular system of
parameters, the ideal is
$I\otimes_{\CC}\CC[x_{5},\dots,x_{n}]$ with $I$ as in \Cref{lem:localproducts},
and \eqref{eq:kunnethext} with the third factor gives
$\Ext^{2}(I_{p},I_{p})=N_{1}\otimes N_{2}\otimes\cO_{W,p}$, where $W$ is the
intersection of the branches. The jet classes along $\partial_{5},\dots,\partial_{n}$
vanish, since these derivations act on a free factor, and the products of the
classes along $\partial_{1},\dots,\partial_{4}$ are as in the lemma. Reducing
modulo the maximal ideal of $\cO_{W,p}$ shows that $a_{\theta}a_{\theta'}\ne0$
whenever $\theta(p)\wedge\theta'(p)$ has a nonzero component in
$N_{1}\otimes N_{2}$. When $n=4$ and $F_{p}$ is the complex
$[\cO\to\cO_{P_{1}}\oplus\cO_{P_{2}}]$, \Cref{lem:localproducts}(ii) applies
directly. The rest of the argument is that of
\Cref{thm:candidatefails}, using \Cref{thm:factor} for $n\ge3$. For the last
assertion, $Z_{i}\cap Z_{j}=D_{i}\cap D_{i+1}\cap D_{j}\cap D_{j+1}$ is a
transversal intersection of four divisors, of codimension four, and at a
general point of it no fifth divisor passes, so $Z$ has exactly the two
branches $Z_{i}$ and $Z_{j}$ there.
\end{proof}

\begin{remark}[Scope of the local obstruction]\label{rem:candidatemeaning}
\Cref{thm:candidatefails} answers the question posed in
\cite[Section 12]{Mar25b}: the secant objects proposed there do not satisfy
the weakened criterion of \cite[Question 11.4]{Mar25b}, for any
discriminant. The reason is local and geometric. That question asks whether
the semiregularity map may be required to be injective only on the image of
the evaluation map. For a secant object this image contains, by
\Cref{thm:factor}, the products $A_{\theta}A_{\theta'}$ of the translation
classes shifted by the scalars
$\tfrac d2(\theta\wedge\theta')\lrcorner\,\Theta^{2}$, and the
semiregularity map kills them all. For the criterion to hold these shifted
products must vanish, and they cannot: at an isolated double point of the
support, the translation classes of the two branches multiply to a nonzero
local class, while a scalar has no local class at all. The fourth ingredient
of \Cref{rem:deformation} can accordingly be stated more sharply. A secant
object of positive rank on a fourfold that satisfies the weakened criterion
must have, at every point, germs whose translation classes multiply to zero
in the local $\Ext^{2}$. By \Cref{lem:localproducts} this excludes every
support with two transversal branches of codimension two through a point,
whether or not the branches are separated, and by
\Cref{cor:localobstruction} it excludes the chain of consecutive
intersections of translates in every dimension. An object supported on a
smooth surface, or on surfaces meeting only along curves, is not excluded by
this argument, but no such object with Chern character on the secant plane
is known.
\end{remark}
